\section{\ours\ dataset} \label{sec:dataset}

The \ours\ dataset consists of \numclaims claims derived from \numpapers biomedical articles, along with supporting evidence, forming \todo expert-annotated claim / label / evidence sets. Sources of claim and evidence are articles from \todo journals. For evidence documents, we hand-selected a small number of high-quality journals in fields broadly considered as biomedicine. The list of journals is shown in Appendix \ref{sec:evidence_journals}. The journals range in focus from basic science (e.g. \emph{Cell}) to clinical trial reports (e.g. the \emph{BMJ}). \todoit{Mention something about structured vs. unstructured abstracts.}

\subsection{Annotators}

Since both claim-writing and evidence-finding tasks require considerable expertise, we recruited undergraduate and graduate students from life sciences departments at the University of Washington to serve as annotators. The claim-writing and evidence-finding stages of the task were performed separately to minimize cognitive load on annotators.

Annotator training took place in two stages. First, the annotators attended an in-person training session for claim writing. Then, they continued to write claims remotely, but were monitored by supervisors and asked to re-write claims that did not satisfy the goals of the task. Some claim writers, other recruited annotators, along with three scientists from the Allen Institute were then trained to provide evidence judgments. Evidence annotators were assigned claims written by someone other than themselves. \todoit{is true?} Feedback on evidence annotations was provided by supervisors, and annotators were given the opportunity to correct annotations.

\subsection{Claim annotation}

\todoit{how citation sentences are selected for rewriting}

Given a citing sentence, annotators are asked to rewrite the sentence as check-able claims. The annotator is provided the context around the citing sentence, along with the title and abstract of the source scientific article. We ask the annotator to write claims based on the primary findings discussed in the citing sentence. Each claim should be about one finding on one aspect of an entity or process. When compound statements are found in the citing sentence, annotators are asked to split the citing sentence into any number of independent claims. They were asked to skip claims that are not checkable, such as \todoit{show an example}. Variety among claim language is encouraged. For example, given a citing sentence:

\todoit{example citation sentence}

\noindent an annotator may rewrite this as \todo or \todo.

A total of \todo annotators -- \todo undergraduate and \todo graduate students -- wrote claims. The most claims written by a single annotator were \todo. In initial experiments, annotators exhibited high agreement (\todo in deciding which claims to skip and which to re-write into claims.)

\subsection{Evidence annotation}

\kyle{here}

Different annotators are then tasked to determine the availability of evidence for each claim. Given a claim and the abstract of the cited article, annotators are asked to select all \textbf{evidence sets} from the abstract text. Each evidence set is a collection of one or more sentences that provide evidence in support of or contradicting the claim. The level of support is represented by a \textbf{stance label}, selected from among the options: \textsc{FullSupport}, \textsc{PartialSupport}, \textsc{PartialRefute}, and \textsc{FullRefute}. Multiple evidence sets can exist for a given claim and they do not have to share the same label. If the abstract does not provide evidence supporting or refuting the claim, annotators select \textsc{NoEvidence}.

As in \citet{Hanselowski2019ARA}, we provide more granular support / contradiction labels. These granular labels allow us to address some of the challenges of scientific fact checking described in Section \ref{sec:challenges}, such as identifying cases of missing qualifiers and/or mechanistic detail.

% \dw{Note that we're not the first ones to do partial support / contradiction; \cite{Hanselowski2019ARA} use a scale of false, mostly false, other, mostly true, true.}

Each sentence in an evidence set is then assigned a \textbf{discourse role}: \textsc{Primary} or \textsc{Background}.  A \textsc{Primary} sentence is required to derive the stance label of the evidence set -- that is, if removed, the evidence set no longer provides support or contradiction for the main body of the claim. Each evidence set must therefore contain at least one \textsc{Primary} sentence. Sentences that are not selected into any evidence set are considered \textsc{Inconclusive}.

\todoit{Example here}.

If an evidence set offers partial support or refutation towards the claim, we ask annotators to provide a rationale to the incompleteness in the form of a \textbf{claim edit}. If \textsc{PartialSupport} is selected, we ask for the claim to be edited (with minimal additions and deletions) until the selected evidence set fully supports the edited claim. If \textsc{PartialRefute} is selected, we correspondingly ask for the claim to be edited until the selected evidence set offers full refutation. For example, \todoit{add example}.

\todoit{merge this example into section \ref{sec:challenges}}
We also ask annotators to consider the specificity of support provided by the evidence. Consider the claim \textit{Smoking causes cancer}. The following annotations show evidence that is more specific and more general to the claim, and the appropriate level of support offered.

\begin{itemize}
    \item Evidence more specific than claim $\to$ \mbox{\textsc{FullSupport}}
    \begin{itemize}
        \item \textit{Smoking cigarettes} causes cancer
        \item Smoking \textit{accelerates} cancer
        \item Smoking causes \textit{lung cancer}
    \end{itemize}

    \item Evidence more general than claim $\to$ \mbox{\textsc{PartialSupport}}
    \begin{itemize}
        \item \textit{Drugs} cause cancer
        \item Smoking \textit{is associated with} cancer
        \item Smoking causes \textit{disease}
    \end{itemize}
\end{itemize}
\kyle{Check whether this is in \fever}
% This asymmetry requires a fact checker to understand general-specific relationships between entities or relations in the claim and evidence.

A total of \todo undergraduates, \todo graduate students, and three scientists at the Allen Institute annotated evidence. Of the student annotators, \todo had previously worked as claim writers. Annotators were shown a claim, and an abstract, and asked to identify all evidence sets that provide support or refutation for the claim \todoit{Show a picture of the interface}.

To measure inter-annotator agreement, \todo claim / evidence pairs were annotated by two annotators. We compute Cohen's $\kappa$ for coarse-grained and fine-grained label agreement, shown in Table \todo. We also compute sentence-level Cohen's $\kappa$ for evidence agreement, as shown in Table \todo. Finally, for comparison with retrieval models, we compute precision, recall, and F1 by  treating one  annotator's evidence as ``gold'' and the other's as ``predicted''. These results are shown in Table \todo.


\begin{itemize}
    \item \kyle{how many claims have  0, 1, 2, 3 evidence sets?}
    \item \kyle{On average, what percentage of the abstract is included in an evidence set (and broken into \textsc{Primary} and \textsc{Background}).}
    \item All unselected sentences are considered \textsc{Inconclusive}. \kyle{Can we justify why didnt separate into RelatedButInconclusive vs Unrelated?  I think compelling argument is the full abstract is almost surely related}
\end{itemize}

% \todoit{move up a section}
% \paragraph{Acronyms, aliases, and synonyms}   Acronyms (e.g. ``Human immunodeficiency virus'' and ``HIV''), aliases (e.g. ``Trastuzumab'' and ``Herceptin''), and synonyms (e.g. ``\kyle{example}'') do not warrant a \textsc{PartialSupport}.
% \kyle{Can we count how many times this occurs?}


\subsection{Dataset statistics}

Figure \todo shows the distribution of abstract lengths (in sentences) for the evidence documents. Statistics for structured and unstructured abstracts are shown separately. Figure \todo shows the distribution of stance labels. Figure \todo shows the distribution of the number of evidence sets per claim, and Figure \todo shows the distribution of the number of sentences per evidence set. The majority of supported claims are supported by a single evidence set, which usually consists of a single primary sentence.